\section{Subgrupos, generación y grupos cíclicos}\label{sec:grupos-subgrupos}

\subsection{Criterio y construcciones de subgrupos}\label{tema:4-2-1}
\begin{defi}
Un grupo $H$ cuyo conjunto subyacente está contenido en $G$ es un \emph{subgrupo} de $G$, escrito $H\le G$, si la inclusión $i:H\hookrightarrow G$ es un homomorfismo.
\end{defi}
En particular, si inicialmente escribimos $\bullet$ para la operación de $H$ y $\cdot$ para la de $G$, entonces
\[
a\bullet b=i(a\bullet b)=i(a)\cdot i(b)=a\cdot b.
\]
La operación de un subgrupo es siempre la inducida por el grupo ambiente.

\begin{prop}[Criterio de subgrupo]\label{prop:C4-criterio}
Un subconjunto no vacío $H\subseteq G$ es un subgrupo si y sólo si $ab^{-1}\in H$ para todos $a,b\in H$.
\end{prop}
\begin{proof}
Si $H\le G$, contiene los inversos y sus productos, de modo que contiene $ab^{-1}$. Recíprocamente, elijamos $a\in H$. Entonces $e=aa^{-1}\in H$; para cualquier $b\in H$ se sigue que $b^{-1}=eb^{-1}\in H$. Por tanto, para $a,b\in H$ también
\[
ab=a(b^{-1})^{-1}\in H.
\]
La asociatividad se hereda de $G$, y el neutro y los inversos recién encontrados cumplen los axiomas de grupo. La inclusión conserva la operación.
\end{proof}

\begin{lema}\label{lema:C4-interseccion}
La intersección de una familia de subgrupos de $G$ es un subgrupo de $G$.
\end{lema}
\begin{proof}
Si $H=\bigcap_{\alpha\in I}H_\alpha$, entonces $e\in H$, pues pertenece a cada $H_\alpha$. Para $a,b\in H$, el elemento $ab^{-1}$ pertenece a cada $H_\alpha$ y, por tanto, a $H$. Aplicamos el criterio anterior. Si $I=\varnothing$, la intersección dentro de $G$ se interpreta como $G$.
\end{proof}

\begin{lema}\label{lema:C4-preimagen}
Si $\varphi:G\to G'$ es un homomorfismo y $H'\le G'$, entonces $\varphi^{-1}(H')\le G$.
\end{lema}
\begin{proof}
Como $\varphi(e_G)=e_{G'}\in H'$, la preimagen es no vacía. Para $a,b$ en ella,
\[
\varphi(ab^{-1})=\varphi(a)\varphi(b)^{-1}\in H'.
\]
Así, $ab^{-1}$ vuelve a estar en la preimagen y se aplica el criterio.
\end{proof}

\subsection{Núcleo e imagen}\label{tema:4-2-2}
\begin{defi}
Para un homomorfismo $\varphi:G\to G'$, definimos su \emph{núcleo} y su \emph{imagen} por
\[
\ker\varphi=\varphi^{-1}(\{e_{G'}\}),\qquad
\im\varphi=\varphi(G).
\]
\end{defi}
El núcleo es un subgrupo por el lema de la preimagen. La imagen también es un subgrupo: contiene $e_{G'}$, y si $a=\varphi(x)$ y $b=\varphi(y)$, entonces
\[
ab^{-1}=\varphi(xy^{-1})\in\im\varphi.
\]

\begin{prop}[Propiedad universal del núcleo]\label{prop:C4-nucleo-universal}
Sea $i:\ker\varphi\hookrightarrow G$ la inclusión. Todo homomorfismo $\alpha:K\to G$ tal que $\varphi\circ\alpha=0_{K,G'}$ se factoriza de manera única como $\alpha=i\circ\bar\alpha$, con $\bar\alpha:K\to\ker\varphi$ un homomorfismo.
\[
\begin{tikzcd}[column sep=large,row sep=large]
K\arrow[r,"\alpha"]\arrow[rr,bend left=25,"0_{K,G'}"]\arrow[dr,destacado,"\exists!\,\bar\alpha"']&G\arrow[r,"\varphi"]&G'\\
&\ker\varphi\arrow[u,hook,"i"']&
\end{tikzcd}
\]
\end{prop}
\begin{proof}
Para todo $k\in K$, $\varphi(\alpha(k))=e_{G'}$, de modo que $\alpha(k)\in\ker\varphi$. Podemos definir $\bar\alpha(k)=\alpha(k)$ con este codominio. Conserva productos porque $\alpha$ los conserva y la operación del núcleo es la restringida de $G$. La igualdad $i\circ\bar\alpha=\alpha$ obliga a esta definición y prueba la unicidad.
\end{proof}
Equivalentemente, $(\ker\varphi,i)$ es terminal en la categoría cuyos objetos son los homomorfismos $\alpha:K\to G$ anulados por $\varphi$, y cuyas flechas $u:(K,\alpha)\to(L,\beta)$ satisfacen $\beta\circ u=\alpha$.

\subsection{Subgrupo generado por un conjunto}\label{tema:4-2-3}
Sea $A\subseteq G$. La inclusión de conjuntos se extiende mediante el grupo libre a un único homomorfismo
\[
\begin{tikzcd}[column sep=large,row sep=large]
F(A)\arrow[r,destacado,"\varphi_A"]&G\\
A\arrow[u,"j"]\arrow[ur,hook]&
\end{tikzcd}
\]
Su imagen es el \emph{subgrupo generado por $A$}, denotado por $\ang{A}$. Mediante palabras,
\[
\ang{A}=\set{a_1^{\varepsilon_1}\cdots a_m^{\varepsilon_m}
\given m\ge0,\ a_i\in A,\ \varepsilon_i\in\{1,-1\}}.
\]
El producto vacío es $e_G$. Esto incluye el caso $A=\varnothing$, para el cual $\ang{A}=\{e_G\}$.

\begin{prop}\label{prop:C4-generado}
El subgrupo $\ang{A}$ es el menor subgrupo de $G$ que contiene a $A$; en concreto,
\[
\ang{A}=\bigcap_{\substack{H\le G\\A\subseteq H}}H.
\]
\end{prop}
\begin{proof}
La imagen de $\varphi_A$ es un subgrupo y contiene cada $a\in A$. Si $H\le G$ contiene a $A$, contiene sus inversos y todos los productos finitos de éstos; por la descripción anterior, $\ang{A}\subseteq H$. Así está contenido en la intersección. A su vez, como $\ang{A}$ es uno de los subgrupos que se intersectan, la intersección está contenida en él.
\end{proof}

Si $G$ es abeliano, en notación aditiva la extensión pasa por $F^{\mathrm{ab}}(A)$ y obtenemos
\[
\ang{A}=\set*{\sum_{i=1}^m n_i a_i\given
m\ge0,\ a_i\in A,\ n_i\in\mathbb Z}.
\]
Si $A=\{a\}$, el grupo libre $F(A)$ es isomorfo a $\mathbb Z$, y $\varphi_A$ se identifica con el homomorfismo exponencial $m\mapsto a^m$. Así,
\[
\ang{a}=\{\ldots,a^{-2},a^{-1},e,a,a^2,\ldots\}.
\]
Un grupo es \emph{finitamente generado} si $G=\ang{A}$ para algún subconjunto finito $A\subseteq G$.

\subsection{Subgrupos de los enteros y de sus cocientes}\label{tema:4-3-1}
Para $d\in\mathbb Z$ escribimos $d\mathbb Z=\ang{d}=\set{m\in\mathbb Z\given d\mid m}$.

\begin{prop}\label{prop:C4-subgrupos-Z}
Todo subgrupo $H\le\mathbb Z$ es $d\mathbb Z$ para un único entero $d\ge0$.
\end{prop}
\begin{proof}
Si $H=\{0\}$, tomamos $d=0$. En otro caso, como $H$ contiene los opuestos de sus elementos, contiene algún entero positivo. Sea $d$ el menor entero positivo de $H$. Todos sus múltiplos pertenecen a $H$, por lo que $d\mathbb Z\subseteq H$. Si $h\in H$, el algoritmo de la división da
\[
h=qd+r,\qquad 0\le r<d.
\]
Pero $r=h-qd\in H$; la minimalidad de $d$ fuerza $r=0$, de modo que $h\in d\mathbb Z$. La unicidad se sigue de que, para $d>0$, éste es el menor elemento positivo de $d\mathbb Z$.
\end{proof}

\begin{prop}\label{prop:C4-subgrupos-Zn}
Si $n>0$ y $H\le\mathbb Z/n\mathbb Z$, existe un único divisor positivo $d$ de $n$ tal que $H=\ang{[d]_n}$.
\end{prop}
\begin{proof}
La preimagen de $H$ bajo $\pi_n:\mathbb Z\to\mathbb Z/n\mathbb Z$ es un subgrupo de $\mathbb Z$, luego es $d\mathbb Z$ para algún $d\ge0$:
\[
\begin{tikzcd}[column sep=large,row sep=large]
\mathbb Z\arrow[r,two heads,"\pi_n"]&\mathbb Z/n\mathbb Z\\
\pi_n^{-1}(H)=d\mathbb Z\arrow[u,hook]\arrow[r,two heads,"\restr{\pi_n}{d\mathbb Z}"']&H\arrow[u,hook]
\end{tikzcd}
\]
Como $n\mathbb Z\subseteq\pi_n^{-1}(H)$, tenemos $n\in d\mathbb Z$, de donde $d>0$ y $d\mid n$. La sobreyectividad de $\pi_n$ implica
\[
H=\pi_n(d\mathbb Z)=\ang{[d]_n}.
\]
Recíprocamente, si $d>0$ divide a $n$, entonces $\pi_n^{-1}(\ang{[d]_n})=d\mathbb Z$: la congruencia $x\equiv kd\pmod n$ implica que $x$ es múltiplo de $d$. La unicidad se reduce así a la del subgrupo de enteros.
\end{proof}
Como todo grupo cíclico es isomorfo a $\mathbb Z$ o a $\mathbb Z/n\mathbb Z$, sus subgrupos son también cíclicos.

\begin{ejem}
La tabla muestra los seis subgrupos de $\mathbb Z/12\mathbb Z$, escribiendo $[a]$ en lugar de $[a]_{12}$.
\begingroup
\renewcommand{\arraystretch}{1.2}
\[
\begin{array}{c|l}
d&\ang{[d]}\\[2pt]\hline
1&\{[0],[1],[2],[3],[4],[5],[6],[7],[8],[9],[10],[11]\}\\
2&\{[0],[2],[4],[6],[8],[10]\}\\
3&\{[0],[3],[6],[9]\}\\
4&\{[0],[4],[8]\}\\
6&\{[0],[6]\}\\
12&\{[0]\}
\end{array}
\]
\endgroup
La correspondencia invierte divisibilidad e inclusión:
\[
d\mid d'\quad\Longleftrightarrow\quad
\ang{[d']_{12}}\subseteq\ang{[d]_{12}}.
\]
El esquema de la \cref{fig:subgrupos-z12} coloca los subgrupos mayores arriba. Las líneas indican inclusiones inmediatas; en el lado de los divisores, el número superior divide al inferior.
\figuraApuntes{Figuras/SubgruposZ12}{fig:subgrupos-z12}{Subgrupos de $\mathbb Z/12\mathbb Z$ y divisores de $12$.}
\end{ejem}

\subsection{Caracterización de los monomorfismos}\label{tema:4-3-2}
\begin{prop}\label{prop:C4-mono}
Para un homomorfismo $\varphi:G\to G'$, son equivalentes:
\begin{enumerate}[label=\textup{(\alph*)}]
\item $\varphi$ es un monomorfismo en $\Grp$;
\item $\ker\varphi=\{e_G\}$;
\item $\varphi$ es inyectivo como función.
\end{enumerate}
\end{prop}
\begin{proof}
Si $\varphi$ es mono, consideremos la inclusión $i:\ker\varphi\to G$ y el homomorfismo trivial $0:\ker\varphi\to G$. Ambos cumplen $\varphi\circ i=\varphi\circ0$. La cancelación da $i=0$, de modo que todo elemento del núcleo es $e_G$.

Si el núcleo es trivial y $\varphi(g_1)=\varphi(g_2)$, entonces
\[
\varphi(g_1g_2^{-1})=\varphi(g_1)\varphi(g_2)^{-1}=e_{G'}.
\]
Por tanto $g_1g_2^{-1}=e_G$ y $g_1=g_2$: la función es inyectiva.

Finalmente, si $\varphi$ es inyectivo y $\alpha,\beta:K\to G$ satisfacen $\varphi\circ\alpha=\varphi\circ\beta$, para cada $k\in K$ se tiene $\varphi(\alpha(k))=\varphi(\beta(k))$. La inyectividad da $\alpha(k)=\beta(k)$ para todo $k$. Por tanto, $\alpha=\beta$, lo que prueba la cancelación requerida.
\end{proof}
